\src{170}
\begin{equation}\tag{1.3}
  H^{p}(X,B)=(R^{p}\Gamma)(B)=H^{p}\bigl(\Gamma(C(B))\bigr).
\end{equation}
Clearly $h_A=\Gamma\hsh_A$; on the other hand, one verifies that $\hsh_A$
takes injective objects to $\Gamma$-acyclic objects. The familiar conclusion is
the following.

\result{PROPOSITION}{1} For every $\Osh$-Module $A$, there exists a
cohomological spectral functor on $\Ccat^{\Osh}$ converging to the graded functor
$\bigl(\operatorname{Ext}_{\Osh}^{*}(X;A,B)\bigr)$, whose initial term is
\begin{equation}\tag{1.4}
  E_{2}^{p,q}(A,B)=H^{p}\bigl(X,\Extsh_{\Osh}^{q}(A,B)\bigr).
\end{equation}
From it one obtains edge homomorphisms and a five-term exact sequence, which we
shall not write down.

\result{COROLLARY}{1} If $A$ is locally isomorphic to $\Osh^{n}$, then there are
canonical isomorphisms
\begin{equation}\tag{1.5}
  \operatorname{Ext}_{\Osh}^{p}(X;A,B)
  \overset{\sim}{\longleftarrow}
  H^{p}\bigl(X,\Homsh_{\Osh}(A,B)\bigr)
\end{equation}
(given by edge homomorphisms of the spectral sequence). In particular, there is
a canonical isomorphism
\begin{equation}\tag{1.6}
  \operatorname{Ext}_{\Osh}^{p}(X;\Osh,B)=H^{p}(X,B).
\end{equation}

To use these results, one must be able to describe the
$\Extsh_{\Osh}^{p}(A,B)$. These functors are computed locally: if $U$ is an open
subset of $X$, then
\[
  \Extsh_{\Osh}^{p}(A,B)|U
  =\Extsh_{\Osh|U}^{p}(A|U,B|U),
\]
as follows from the fact that the restriction to $U$ of an injective
$\Osh$-Module is an injective $(\Osh|U)$-Module. Furthermore, for fixed
$x\in X$, there are functorial homomorphisms
\begin{equation}\tag{1.7}
  \Homsh_{\Osh}(A,B)_x
  \longrightarrow
  \operatorname{Hom}_{\Osh_x}(A_x,B_x),
\end{equation}
which extend uniquely to a homomorphism of cohomological functors in $B$:
\begin{equation}\tag{1.8}
  \Extsh_{\Osh}^{p}(A,B)_x
  \longrightarrow
  \operatorname{Ext}_{\Osh_x}^{p}(A_x,B_x).
\end{equation}

\result{PROPOSITION}{2} If, in a neighborhood of $x$, $A$ is isomorphic to the
cokernel of a homomorphism $\Osh^{m}\longrightarrow\Osh^{n}$, then (1.8) is an isomorphism in degree $p$ when $A$ admits locally a finite free partial resolution through degree $p+1$; it is an isomorphism for every $p$ when $A$ is pseudo-coherent. A single finite presentation alone gives only the degree-zero comparison.\footnote{The printed text reads “(1.7)”; the presence
of $p$ and the preceding sentence require “(1.8).”} This holds in particular if
$A$ is a coherent $\Osh$-Module [3].

